\chapter{Análise do Consumo de Energia Elétrica}



Esta seção apresenta a caracterização do fornecimento de energia elétrica e o perfil de consumo da unidade consumidora avaliada, com base nos dados extraídos das últimas 12 faturas disponibilizadas, complementados por informações técnicas levantadas durante a inspeção de campo.

A análise tem como objetivo identificar o comportamento da demanda e do consumo de energia elétrica, bem como compreender a distribuição do consumo entre os principais usos finais, subsidiando a proposição de medidas de eficiência energética.

\section{Características do Fornecimento}

A Tabela~\ref{tab:caracteristicas-fornecimento} sintetiza os principais parâmetros associados ao ponto de fornecimento da unidade consumidora, incluindo modalidade tarifária, nível de tensão, demanda contratada e outros dados cadastrais relevantes. Estas informações são fundamentais para a compreensão do enquadramento tarifário e do potencial de economia com base em ações técnico-institucionais.


\begin{table}[h]
\centering
\begin{tabular}{|>{\columncolor[HTML]{1a3a5c}}p{8cm}|p{7cm}|}
\hline
\textcolor{white}{\textbf{Grupo de Tensão}} & <<grupoTensao>> \\
\hline
\textcolor{white}{\textbf{Modalidade tarifária}} & <<modalidadeTarifaria>> \\
\hline
\textcolor{white}{\textbf{Demanda contratada}} & <<demandaContratada>> kW \\
\hline
\textcolor{white}{\textbf{Número da instalação na distribuidora}} & <<codInstalacao>> \\
\hline
\textcolor{white}{\textbf{Código da unidade consumidora}} & <<codigoUc>> \\
\hline
\textcolor{white}{\textbf{Classe de consumo}} & <<classeConsumo>> \\
\hline
\end{tabular}
\caption{Características do fornecimento de energia elétrica}
\label{tab:caracteristicas-fornecimento}
\end{table}




%\begin{table}[H]
 %   \centering
  %  \includegraphics[width=\textwidth]{Figuras/Tabela 3.png}
   % \caption{Características do fornecimento de energia elétrica}
%\label{tab:caracteristicas-fornecimento}
%\end{table}





\section{Distribuição do Consumo por Uso Final}

Com base nas observações pontuais, no inventário técnico e nas informações operacionais obtidas durante a visita de campo, foi estimada a participação percentual dos principais usos finais no consumo total de energia elétrica da unidade. Essa análise permite identificar os focos predominantes de consumo e orientar a priorização das medidas de eficiência energética.

Os sistemas de \textbf{iluminação}, \textbf{climatização}, \textbf{sistemas motrizes}, \textbf{refrigeração} e \textbf{outros usos diversos} representam as principais usos finais analisados. A Figura~\ref{fig:distribuicao-uso-final} apresenta a distribuição estimada de consumo por uso final, permitindo visualizar a contribuição relativa de cada categoria e subsidiar decisões técnicas quanto à implementação das ações propostas.

\vspace{1cm}

\begin{figure}[H]
\centering
\includegraphics[width=0.9\textwidth]{Figuras/Grafico 1.png}
\caption{Distribuição estimada do consumo de energia elétrica por uso final}
\label{fig:distribuicao-uso-final}
\end{figure}